\section{Conditioning regimes from objective geometry}
\label{sec:classifications}

The diameter law separates two questions: how quickly the objective
sublevel sets shrink, and how their shape determines a barrier Hessian.
Error bounds answer the first. Their connection to weak sharp minima and
conic degeneracy is classical \citep{burkeferris1993,sturm2000,lourenco2021}.
Here every two-sided estimate of $D(g)$ gives a two-sided conditioning
rate for every fixed self-concordant barrier.

\begin{corollary}[Error bounds and attained rates]
\label{cor:error-dictionary}
Suppose, for some $M>0$, $\theta\in(0,1]$, and all sufficiently small
positive gaps,
\begin{equation}
 \operatorname{dist}(y,\Phi)\leq M g(y)^\theta.
 \label{eq:error-bound}
\end{equation}
Then
\begin{equation}
 \max\left\{\diam\Phi,\frac{g}{\norm{c_{\V}}}\right\}
 \leq D(g)\leq\diam\Phi+2Mg^\theta.
 \label{eq:error-diameter}
\end{equation}
If $\Phi$ is a singleton, every fixed barrier consequently has
$\widehat\kappa_F(g)=O(g^{2\theta-2})$.
More generally, whenever $D(g)=\Theta(g^\theta)$ for
$\theta\in[0,1]$,
\begin{equation}
 \widehat\kappa_F(g)=\Theta(g^{2\theta-2}).
 \label{eq:attained-dictionary}
\end{equation}
For a fixed barrier, the same rates hold with $\mu$ in place of $g$.
The gap statements also hold for barriers with uniformly bounded
parameters, at points with Newton decrements at most a fixed $\bar\rho<1$.
\end{corollary}

\begin{proof}
For two points of $L(g)$, choose nearest points in the compact set
$\Phi$ and apply the triangle inequality. Points already in $\Phi$ have zero
distance, so this proves the upper bound for small $g$.
The inclusion $\Phi\subset L(g)$ proves its first lower bound.
For every $0<g<\Delta$ there is a point $y\in P$ with gap exactly $g$,
by taking a segment from an optimum to a maximizer. Then
$g\leq\norm{c_{\V}}\norm{y-x^*}$ for any $x^*\in\Phi$,
which proves the second lower bound. Apply
\cref{thm:diameter-law,prop:gap-parametrization} to finish.
\end{proof}

At a unique optimum, an error bound alone bounds the conditioning rate
only from above. A matching lower bound requires sublevel chords
whose lengths have the same exponent. This distinction matters for
singularity-degree bounds, as the paired examples in
\cref{sec:fractional-sdp} show.

\begin{proposition}[Faces, vertices, and curved optima]
\label{prop:three-regimes}
For every fixed barrier, the following statements hold.
\begin{enumerate}
 \item If $\dim\Phi>0$, then
 $D(g)=\Theta(1)$ and $\widehat\kappa_F(g)=\Theta(g^{-2})$.
 \item If $P$ is a polytope and $\Phi=\{x^*\}$, then
 $D(g)=\Theta(g)$ and $\widehat\kappa_F(g)=\Theta(1)$.
 No nondegeneracy assumption on the vertex is needed.
 \item If $\Phi=\{x^*\}$ and $D(g)=\Theta(\sqrt g)$, then
 $\widehat\kappa_F(g)=\Theta(g^{-1})$.
\end{enumerate}
\end{proposition}

\begin{proof}
A positive-dimensional optimal face has positive diameter, while
$D(g)\leq\diam P$. For the polytope assertion let $v_1,\ldots,v_k$
be its vertices other than $x^*$ and put
$\delta=\min_i g(v_i)>0$. Every point can be written
$y=\alpha_0x^*+\sum_i\alpha_iv_i$, with nonnegative coefficients
summing to one. Thus
\[
 \norm{y-x^*}\leq\diam(P)\sum_i\alpha_i
 \leq\frac{\diam(P)}{\delta}g(y).
\]
This gives $D(g)=O(g)$; the lower bound in
\eqref{eq:error-diameter} gives the converse. All three conditioning
claims now follow from \eqref{eq:attained-dictionary}.
\end{proof}

These three cases are examples, not a classification of convex sets.
The diameter law applies even when $D(g)$ does not behave like a power
of $g$. For a fixed compact LP, however, the optimum is a vertex or a
positive-dimensional face, so the asymptotic conditioning exponent is
zero or two. An intermediate slope of $\log\widehat\kappa_F$ against
$\log g$ over a finite range of gaps therefore does not establish a
fractional asymptotic rate.

The polytope argument is the elementary finite-vertex form of weak
sharpness. It also shows that primal degeneracy alone cannot make the
condition number of the reduced Hessian unbounded as
$g\downarrow0$. The bound is per instance; its constant need not
be small or uniform over the data.
\Cref{sec:lp-limits} computes the bounded limit for the logarithmic
barrier.

\begin{example}[A curved singleton in a $2\times2$ spectrahedron]
\label{ex:curved-singleton}
Minimize $X_{22}$ over $X\succeq0$, $\tr X=1$. Write
\[
 X=\begin{pmatrix}1-u&v\\v&u\end{pmatrix},\qquad
 0\leq u\leq1,\qquad v^2\leq u(1-u).
\]
The unique optimum is $\Diag(1,0)$. If $u\leq g$, the Frobenius
distance to this optimum is $O(\sqrt g)$. The two feasible points with
$u=g$ and $v=\pm\sqrt{g(1-g)}$ are separated by
$2\sqrt{2g(1-g)}$, so $D(g)=\Theta(\sqrt g)$. Hence every fixed
barrier has condition number $\Theta(g^{-1})$.

For $F=-\log\det X$, symmetry forces $v=0$ on the central path.
Centrality in $u$ gives
\[
 \mu=\frac{g(1-g)}{1-2g},\qquad 0<g<\tfrac12.
\]
The coordinate directions $\partial X/\partial u$ and
$\partial X/\partial v$ are orthogonal and both have squared Frobenius
norm two. The two eigenvalues of the tangent Hessian are therefore
\[
 \frac{1}{g(1-g)},\qquad
 \frac12\bigl(g^{-2}+(1-g)^{-2}\bigr).
\]
They are in increasing order for small $g$, and
$\widehat\kappa_{\log}(g)\sim(2g)^{-1}$.
\end{example}

\subsection{A uniform plateau near a nearly optimal edge}

Consider the fixed simplex with a varying objective,
\begin{equation}
 P=\{x\in\R^3:x_1+x_2+x_3=1,\ x\geq0\},\qquad
 c_\vartheta=(0,\vartheta,1),\qquad 0<\vartheta\leq1.
 \label{eq:near-simplex}
\end{equation}
For each positive $\vartheta$ the unique optimum is $e_1$; at
$\vartheta=0$ the optimal set becomes an edge. The family separates the
asymptotic rate, which is bounded, from a large constant caused by the
nearly optimal edge. As $g$ decreases, the condition number grows like
$g^{-2}$ until $g\approx\vartheta$ and then stays of order
$\vartheta^{-2}$, the plateau (\cref{prop:simplex-plateau}).

\begin{proposition}[The two gap scales of the simplex]
\label{prop:simplex-plateau}
Uniformly for $0<\vartheta\leq1$ and $0<g\leq1$,
\begin{equation}
 D_\vartheta(g)=\Theta\bigl(\min\{g/\vartheta,1\}\bigr).
 \label{eq:simplex-diameter}
\end{equation}
For any family of barriers with parameters at most $\bar\nu$, and for
all $g$ below a threshold independent of $\vartheta$,
\begin{equation}
 \widehat\kappa_{F,\vartheta}(g)
 =\Theta\bigl(\min\{\vartheta^{-2},g^{-2}\}\bigr),
 \label{eq:simplex-plateau}
\end{equation}
with constants depending only on $\bar\nu$. The same estimate holds
at approximately central points, where they exist, with a common Newton-decrement
bound $\bar\rho<1$; the constants then also depend on $\bar\rho$.
\end{proposition}

\begin{proof}
For $x\in L_\vartheta(g)$,
$\vartheta(x_2+x_3)\leq\vartheta x_2+x_3\leq g$, while
$x_2+x_3\leq1$. Hence
\[
 \norm{x-e_1}\leq\sqrt2\min\{g/\vartheta,1\}.
\]
The edge point
$(1-t,t,0)$, $t=\min\{g/\vartheta,1\}$, belongs to
$L_\vartheta(g)$ and has distance $\sqrt2t$ from $e_1$.
Taking diameters proves \eqref{eq:simplex-diameter} with absolute
constants. The set $P$ and its interior-ball radius are fixed, its
objective range is $\Delta=1$, and
\[
 \norm{(c_\vartheta)_{\V}}^2
 =\frac23(\vartheta^2-\vartheta+1)\in[1/2,2/3].
\]
The diameter law is therefore uniform in $\vartheta$.
The common attained gap interval follows from
\cref{cor:uniform-rates}.
\end{proof}

For the logarithmic barrier the limiting plateau height is explicit:
\begin{equation}
 \lim_{\mu\downarrow0}\kappa_{\log,\vartheta}(x_{\log,\vartheta}(\mu))
 =\frac{1+\vartheta^2+\sqrt{1-\vartheta^2+\vartheta^4}}
        {1+\vartheta^2-\sqrt{1-\vartheta^2+\vartheta^4}},
 \label{eq:simplex-exact-limit}
\end{equation}
and multiplying this limit by $\vartheta^2$ and then letting
$\vartheta\downarrow0$ gives $4/3$. To prove this, write the equality
multiplier as $\alpha_\mu=\mu/x_1>0$. Centrality gives
\[
 x_2=\frac{\mu}{\vartheta+\alpha_\mu},\qquad
 x_3=\frac{\mu}{1+\alpha_\mu},\qquad
 \alpha_\mu\longrightarrow0.
\]
The last limit follows from $x(\mu)\to e_1$, itself implied by
$g(\mu)\to0$. On the tangent vector $h=(-u-v,u,v)$ the Rayleigh
quotient of $\mu^2H_{\log}$ converges to
\[
 \frac{\vartheta^2u^2+v^2}{2(u^2+uv+v^2)}.
\]
Its eigenvalues in the fixed Euclidean metric are
$(1+\vartheta^2\pm\sqrt{1-\vartheta^2+\vartheta^4})/3$,
which proves \eqref{eq:simplex-exact-limit}. Their product is
$\vartheta^2/3$, and the larger tends to $2/3$ as
$\vartheta\downarrow0$, proving the stated $4/3$ limit.
This is an iterated limit: $\mu$ first tends to zero with
$\vartheta$ fixed. The jointly uniform statement is
\eqref{eq:simplex-plateau}, which describes the transition near
$g\asymp\vartheta$ without interchanging limits.
